% !TEX root = ../../hw01.tex
% Source: ma-gy7033fall2026homework1.pdf, p. 2, Problem 4(a--b).
\begin{exercise}[A parameter that changes the span]
  \label{ex:hw01:parameter-dependent-span}
  Let \(\mathbb F\) denote \(\R\) or \(\C\). For \(a\in\mathbb F\), set
  \[
    v_1=(1,0,1),\qquad v_2=(0,1,1),\qquad v_3=(1,1,a).
  \]
  \begin{enumerate}[(a)][2]
    \item Determine exactly when \(\{v_1,v_2,v_3\}\) is a basis of
      \(\mathbb F^3\). For each remaining value of \(a\), find a basis
      and the dimension of its span, and describe the span by a linear equation.
    \item When the three vectors form a basis, express any
      \((x,y,z)\in\mathbb F^3\) as their linear combination.
  \end{enumerate}
  \solutionlink{hw01:parameter-dependent-span}
\end{exercise}
